\appendix
\section{Search Protocol (PRISMA-Adapted)}
\label{app:search}

\subsection*{Research question}
We aimed to identify papers that import a modern mathematical or machine-learning
framework into granular computing's core pipeline (neighbourhood approximation,
attribute reduction, positive-region dependency), published between 2000 and 2025.
Using PICO framing: \textbf{P}~= granular computing core (approximation space,
granulation, reduct); \textbf{I}~= imported framework (TDA, topology, conformal
prediction, weak supervision); \textbf{C}~= baseline GrC methods (NRS, GBNRS,
discernibility matrix); \textbf{O}~= classification accuracy, computational
efficiency, theoretical bounds, sufficiency/limits.

\subsection*{Sources}
Literature was retrieved from a semantic scholarly search engine and the OpenAlex
bibliographic database; supplementary forward-citation checks were performed
manually on Google Scholar and Semantic Scholar.  Scopus and Web of Science were
not searched in this study.

\subsection*{Queries}
We ran 20 targeted keyword queries (five per thematic axis: granule
geometry, abstract topology, computational TDA, efficiency, sufficiency/limits)
plus 5 citation-graph queries anchored on the five most-cited papers per
axis (Xia 2020, Zhu 2007, Dai 2018, Kindelan 2021, Santos 2022).  Query strings
are listed in the supplementary file \texttt{search\_protocol.md} released with
this paper.

\subsection*{Inclusion and exclusion criteria}
\textit{Included}: peer-reviewed venues $\geq$ Q2 (JCR) or top conference
(AAAI, IJCAI, ICML, NeurIPS, KDD), or citation count $\geq 20$ for older work;
directly relevant to $\geq 1$ thematic axis.  \textit{Excluded}: domain-only
applications without theoretical novelty; non-peer-reviewed workshop papers;
zero-citation papers published before 2023; redundant companion papers (most
recent version retained).

\subsection*{PRISMA flow}

Figure~\ref{fig:prisma} shows the PRISMA-adapted flow diagram for this
survey.

\begin{figure}[pos=!htbp]
\centering
\begin{tikzpicture}[
  box/.style = {draw, rounded corners=4pt, minimum width=4.2cm,
                minimum height=0.9cm, align=center, font=\small,
                fill=gray!8},
  excl/.style = {draw, rounded corners=3pt, minimum width=3.0cm,
                 minimum height=0.7cm, align=center, font=\scriptsize,
                 fill=red!6, draw=red!40!black},
  add/.style  = {draw, rounded corners=3pt, minimum width=3.0cm,
                 minimum height=0.7cm, align=center, font=\scriptsize,
                 fill=green!8, draw=green!50!black},
  arr/.style  = {-{Stealth[length=4pt]}, thick},
]

\node[box] (id) at (0,0) {\textbf{Identification}\\214 records (database search)};
\node[excl, right=0.8cm of id] (dup) {0 duplicates\\removed};
\draw[arr] (id) -- (dup);

\node[box, below=0.9cm of id] (scr) {\textbf{Abstract screening}\\214 records screened};
\node[excl, right=0.8cm of scr] (e1) {155 outside scope\\excluded};
\draw[arr] (id) -- (scr);
\draw[arr] (scr) -- (e1);

\node[box, below=0.9cm of scr] (ft) {\textbf{Full-text assessment}\\59 records assessed};
\node[excl, right=0.8cm of ft] (e2) {11 excluded\\(domain-only, redundant)};
\draw[arr] (scr) -- (ft);
\draw[arr] (ft) -- (e2);

\node[box, below=0.9cm of ft] (main) {\textbf{Main search}\\48 included};
\node[add, right=0.8cm of main] (cit) {+4 citation chain\\(Jensen, Cornelis, Yao $\times$2)};
\draw[arr] (ft) -- (main);
\draw[arr] (cit) -- (main);

\node[box, below=0.9cm of main, fill=blue!8]
  (final) {\textbf{Final included: 52 papers}\\+ 11 background references};
\draw[arr] (main) -- (final);

\end{tikzpicture}
\caption{PRISMA-adapted flow diagram.  Records were retrieved from
scholarly databases; supplementary forward-citation checks on
Google Scholar and Semantic Scholar.  Scopus and Web of Science were not
searched (Section~\ref{sec:limitations}).}
\label{fig:prisma}
\end{figure}

After deduplication, 214 unique records were identified.
Abstract screening removed 155 records outside scope, leaving 59 eligible.
Full-text assessment of these 59 records excluded 11 as outside scope
(domain-only applications, redundant companion papers), yielding 48 records
for inclusion from the main search.  Backward and forward citation-chain
checks on high-citation anchor papers identified 4 additional papers
(Jensen~\&~Shen~2007, Cornelis et~al.~2010, Yao~2010, Yao~2016) that met
all inclusion criteria, bringing the total to 52 included papers.
A further 11 references were retained as background or foundational
(TDA stability theorems, Cover~\&~Hart~1967, survey references).
The final included set of 52 papers is recorded in
\texttt{literature\_log.csv} (released with this paper).

\subsection*{Data extraction}
Each included paper was coded in \texttt{literature\_log.csv} with fields:
paper identifier, first author, year, venue, SJR quartile, citation count,
thematic axis (1/2a/2b/3/4/5), relevance level (core/peripheral/background),
failure-pattern tag (A/B/C/none), read level (abstract/full), key claim,
key finding, and notes.  The \texttt{read\_level} field records the depth of
data extraction (abstract-level coding vs.\ full-text analysis of methods and
results); the ``full-text assessment'' in the PRISMA flow above refers to the
inclusion/exclusion screening step, not the coding depth.  All 52 included
papers were read at full-text depth for the screening decision; data
extraction relied on abstract and methods sections for most papers, with full
experimental re-analysis limited to the ten papers directly involved in the
failure-pattern evidence (Sections~\ref{sec:failure}--\ref{sec:methodology}).
Failure-pattern tagging: A = reduces to known granular
quantity; B = dominance of label-orthogonal geometry; C = self-optimising
evaluation.

\subsection*{Full-text coding of reporting practice}
\label{app:coding}
We coded the reporting practice of 22 papers we were able to read at full
text---13 of them members of the Axis-1 corpus above---released as
\texttt{fulltext\_coding.csv}.  Each row records four presence/absence
fields, judged from the experimental section alone:
\begin{enumerate}[leftmargin=*,itemsep=1pt]
  \item \texttt{all\_features\_baseline} --- does the paper report a
        classification result on the unreduced attribute set, so the reader
        can see whether selection helped?  A mention of the ``original
        dataset'' in a method description does not count; a reported accuracy
        does.
  \item \texttt{reports\_dispersion} --- does at least one results table give
        a standard deviation or per-fold value alongside a mean?
  \item \texttt{significance\_test} --- is a named statistical test
        (Friedman, Wilcoxon, Nemenyi, Bonferroni--Dunn, $t$-test) actually
        carried out on the paper's own results?  Citing a test in related
        work does not count.
  \item \texttt{hyperparam\_ablation} --- does the paper report results at
        more than one value of any hyperparameter of its own method?  Varying
        an external experimental condition, such as an injected noise level,
        does not count.
\end{enumerate}
Every cell carries the quotation or the structural observation it rests on in
the \texttt{evidence} column, so an individual judgement can be checked and
overturned without re-reading the paper.  The fields deliberately record
presence rather than quality: a paper credited with a significance test is
not thereby credited with an appropriate one.  The papers are those we could obtain in full text.
Thirteen of the 22 are members of the Axis-1 corpus above, and constitute 13
of its 25 papers; the remaining nine are from the same literature but outside
our inclusion window.  Section~\ref{sec:significance} therefore reports counts
for both groupings rather than percentages of the field.
Table~\ref{tab:practice_detail} gives the per-paper judgements for the corpus
members.

% AUTO-GENERATED by gen_tables.py -- do not edit by hand.
% Regenerate:  python3 gen_tables.py   Audit:  python3 verify_tables.py
\begin{table}[pos=!htbp]
\caption{Per-paper coding for the 13 Axis-1 corpus papers read at
full text.  B~= all-features baseline; D~= dispersion; S~= significance test;
A~= own-hyperparameter ablation.  $\bullet$~= present, --~= absent.  The
quotation supporting each judgement is in \texttt{fulltext\_coding.csv}.}
\label{tab:practice_detail}
\centering\small
\begin{tabular}{llcccc}
\toprule
Log ID & Paper & B & D & S & A \\
\midrule
T1\_01 & Xia 2020 & -- & -- & -- & -- \\
T1\_02 & Sun 2024 & -- & -- & $\bullet$ & $\bullet$ \\
T1\_03 & Zhang 2025 & $\bullet$ & $\bullet$ & $\bullet$ & -- \\
T1\_04 & Sun 2025 & -- & -- & $\bullet$ & -- \\
T1\_06 & Qian 2024 & -- & -- & $\bullet$ & -- \\
T1\_08 & Ju 2023 & -- & $\bullet$ & -- & -- \\
T1\_09 & Xia 2022 & $\bullet$ & -- & -- & $\bullet$ \\
T1\_10 & Xia 2023 & $\bullet$ & -- & -- & -- \\
T1\_11 & Yang 2024 & -- & -- & $\bullet$ & -- \\
T1\_13 & Xie 2024 & -- & -- & -- & -- \\
T1\_17 & Chen 2021 & -- & -- & $\bullet$ & -- \\
T1\_18 & Peng 2022 & -- & -- & -- & -- \\
T1\_20 & Xia 2024 & -- & -- & $\bullet$ & $\bullet$ \\
\bottomrule
\end{tabular}
\end{table}

